\section{Related Work}\label{sec:related}

\paragraph{Artificial immune systems.}
Hofmeyr and Forrest~\cite{hofmeyr2000architecture} proposed an
architecture for distributed intrusion detection modeled on the
complement cascade and T-cell activation.  Dasgupta~\cite{dasgupta2006advances}
surveys two decades of AIS research.  These systems demonstrate that
immune-inspired pattern matching can detect novel threats, but
comparative benchmarks against non-biological baselines (e.g., rule-based
detectors) are sparse.  Our quorum sensing benchmark addresses a related
problem---multi-agent threat consensus---but explicitly compares the
biological signal-accumulation model against majority voting and
independent detection.

\paragraph{Swarm intelligence.}
Bonabeau, Dorigo, and Theraulaz~\cite{bonabeau1999swarm} established
the foundations of swarm-based optimization (ACO, PSO).  These
algorithms use biological coordination metaphors (pheromone trails,
flocking rules) for optimization.  Operon's quorum sensing addresses a
different problem: \emph{consensus} rather than \emph{optimization}.
The distinction matters: our benchmark tests whether the biological
coordination model achieves better precision--recall tradeoffs than
simpler voting schemes, not whether it finds better optima.

\paragraph{Metabolic computing.}
Resource-aware agent systems have been explored in economic
frameworks~\cite{wellman2001auction}, but explicitly metabolic models---with
multiple energy currencies, state-dependent behavior, and
regeneration---are less common.  Operon's ATP\_Store draws on AMPK
signaling pathway structure (KEGG hsa04152), where the AMP:ATP ratio
and its rate of change drive metabolic state transitions.  Our
benchmark tests whether this biological structure provides measurable
benefits over a flat budget counter.

\paragraph{Free energy principle.}
Friston's free energy principle~\cite{friston2010} provides the
theoretical foundation for Operon's epiplexity monitor: healthy agents
minimize surprise while maintaining viability; agents with high
perplexity and low novelty are in pathological loops.  Our benchmark
tests whether this theoretical framework translates to practical
stagnation detection advantages over cosine-similarity thresholds.

\paragraph{Multi-agent benchmarks.}
Recent work on multi-agent scaling laws~\cite{kim2025scaling} establishes
architecture-level performance predictors.  Operon's epistemic topology
(Paper~1) derives similar bounds from wiring diagram structure.  Our
benchmarks complement this work by testing \emph{mechanism-level}
guarantees rather than \emph{architecture-level} predictions.
